\documentclass[11pt,answers]{exam} %noanswers
\usepackage{multicol}
\usepackage[usenames]{color} %used for font color
\usepackage{amssymb} %maths
\usepackage{amsmath} %maths
\usepackage[utf8]{inputenc} %useful to type directly diacritic characters
\usepackage{siunitx}
\usepackage{mathtools}
\usepackage{float}
\usepackage[version=4]{mhchem}
\newcommand{\qtty}[2]{\left(\qty{#1}{#2}\right)}
\sisetup{display-per-mode = symbol }
\sisetup{inter-unit-product = \ensuremath { { } \cdot { } } }
\usepackage[letterpaper, total={7in, 10in}]{geometry}
\usepackage{graphicx} % Required for inserting images
\renewcommand{\epsilon}{\varepsilon}
\renewcommand{\vec}[1]{\mathbf{#1}}
\newcommand{\mat}[1]{\begin{bmatrix*}[r] #1 \end{bmatrix*}}
\newcommand{\R}[1]{\mathbb{R}^{#1}}

\begin{document}
\flushleft
Name:\ \makebox[4in]{\hrulefill} \hspace{1in} Date:\ \makebox[1in]{\hrulefill}
\begin{center}
    \textbf{\S 2.3 Characteristics of Invertible Matrices}
\end{center}


\begin{questions}
    \uplevel{Unless otherwise specified, assume that all matrices in these exercises are $n\times n$. Determine which of the matrices in Exercises 1–10 are invertible. Use as few calculations as possible. Justify your answers}
    \question $ \mat{5 & 7 \\ -3 & -6} $
    \question $ \mat{-4 & 6 \\ 6 & -9} $
    \question $ \mat{ 5 & 0 & 0 \\ -3 & -7 & 0 \\ 8 & 5 & -1} $
    \question $ \mat{ -7 & 0 & 4 \\ 3 & 0 & -1 \\ 2 & 0 & 9} $
    \question $ \mat{0 & 4 & 7 \\ 1 & 0 & 5 \\ -5 & 8 & -2} $
    \question $ \mat{1 & -5 & -4 \\ 0 & 3 & 4 \\ -3 & 6 & 0} $
    \question $ \mat{-1 & 0 & 2 & 1 \\ -5 & -3 & 9 & 3 \\ 3 & 0 & 1 & -3 \\ 0 & 3 & 1 & 2} $
    \question $ \mat{1 & 3 & 7 & 4 \\ 0 & 5 & 9 & 6 \\ 0 & 0 & 2 & 8 \\ 0 & 0 & 0 & 10} $
    \question $ \mat{4 & 0 & -7 & -7 \\ -6 & 1 & 11 & 9 \\ 7 & -5 & 10 & 19 \\ -1 & 2 & 3 & -1} $
    \question $ \mat{5 & 3 & 1 & 7 & 9 \\ 6 & 4 & 2 & 8 & -8  \\ 7 & 5 & 3 & 10 & 9 \\ 9 & 6 & 4 & -9 & -5 \\ 8 & 5 & 2 & 11 & 4}$
    %=============
    \begin{center}\rule{4in}{0.1mm}\end{center}
    \uplevel{In Exercises 11-20, the matrices are all $n \times n$. Each part of the exercises is an \emph{implication} of the form ``If `statement 1', then `statement 2'." Mark an implication True is the truth of `` statement 2'' \emph{always} follows whenever ``statement 1'' happens to be true. An implication is False if there is an instance in which ``statement 2'' is false but ``statement 1'' is true. Justify each answer.}
    \question (\textbf{T}/\textbf{F}) If the equation $A\vec{x} = \vec{0}$ has only the trivial solution, then $A$ is row equivalent to the $n \times n$ identity matrix.
    \question (\textbf{T}/\textbf{F}) If there is an $n \times n$ matrix $D$ such that $AD = I$, then there is also an $n \times n$ matrix $C$ such that $CA = I$.
    \question (\textbf{T}/\textbf{F}) If the columns of $A$ span $\mathbb{R}^n$, then the columns are linearly independent.
    \question (\textbf{T}/\textbf{F}) If the columns of $A$ are linearly independent, then the columns of $A$ span $\mathbb{R}^n$.
    \question (\textbf{T}/\textbf{F}) If $A$ is an $n \times n$ matrix, then the equation $A\vec{x} = \vec{b}$ has at least one solution for each $\vec{b}$ in $\mathbb{R}^n$.
    \question (\textbf{T}/\textbf{F}) If the equation $A\vec{x} = \vec{b}$ has at least one solution for each $\vec{b}$ in $\mathbb{R}^n$, then the solution is unique for each $\vec{b}$.
    \question (\textbf{T}/\textbf{F}) If the equation $A\vec{x} = \vec{0}$ has a nontrivial solution, then $A$ has fewer than $n$ pivot positions.
    \question (\textbf{T}/\textbf{F}) If the linear transformation $\vec{x} \mapsto A\vec{x}$ maps $\mathbb{R}^n$ to $\mathbb{R}^n$, then $A$ has $n$ pivot positions.
    \question (\textbf{T}/\textbf{F}) If $A^T$ is not invertible, then $A$ is not invertible.
    \question (\textbf{T}/\textbf{F}) If there is a $\vec{b}\in\mathbb{R}^n$ such that the equation $A\vec{x}=\vec{b}$ is inconsistent, then the transformation $\vec{x} \mapsto A\vec{x}$ is not one-to-one.
    %=============
    \begin{center}\rule{4in}{0.1mm}\end{center}
    \question An $m \ times n$ \textbf{upper triangular matrix} is one whose entires \emph{below} the main diagonal are 0's (as in Exercise 8). When is a square upper triangular invertible? Justify your answer.
    \question An $m \ times n$ \textbf{lower triangular matrix} is one whose entires \emph{above} the main diagonal are 0's (as in Exercise 8). When is a square upper triangular invertible? Justify your answer.
    \question Can a square matrix with two identical columns be invertible? Why or why not?
    \question Is it possible for a $5 \times 5$ matrix to be invertible when its columns do not span $\mathbb{R}^5$? Why or why not?
    \question If $A$ is invertible, then the columns of $A^{-1}$ are linearly independent. Explain why.
    \question If $C$ is $6 \times 6$ and the equation $C\vec{x} = \vec{v}$ is consistent for every $\vec{v}$ in $\mathbb{R}^6$, is it possible that for some $\vec{v}$, the equation $C\vec{x}=\vec{v}$ has more than one soltuion? Why or why not?
    \question If the columns of a $7 \times 7$ matrix $D$ are linearly independent, what can you say about the solutions of $D\vec{x}=\vec{b}$? Why?
    \question If $n \times n$ matrices $E$ and $F$ have the property that $EF = I$, then $E$ \& $F$ commute. Explain why.
    \question If the equation $G\vec{x}=\vec{y}$ has more than one solution for some $\vec{y}$ in $\mathbb{R}^n$, can the columns of $G$ span $\mathbb{R}^n$? Why or why not?
    \question If the equation $H\vec{x}=\vec{c}$ is consistent for some $c$ in $\mathbb{R}^n$, what can you say about the equation $H\vec{x}=\vec{0}$? Why?
    \question If an $n \times n$ matrix $K$ cannot be row reduced to $I_n$, what can you say about the columns of $K$? Why?
    \question If $L$ is $n \times n$ and the equation $L\vec{x}=\vec{0}$ has the trivial solutions, do the columns of $L$ span $\mathbb{R}^n$? Why?
    \question Verify the boxed statement preceding Example 1.
    \question Explain why the columns of $A^2$ span $\mathbb{R}^n$ whenever the columns of $A$ are linearly independent.
    \question Show that if $AB$ is invertible, so is $A$. You cannot use Theorem 6(b), because you cannot \emph{assume} that $A$ and $B$ are invertible. Hint: There is a matrix $W$ such that $ABW = I$ Why?
    \question Show that if $AB$ is invertible, so is $B$.
    \question If $A$ is an $n \times n$ matrix and the equation $A\vec{x}=\vec{b}$ has more than one solution for some $\vec{b}$, then the transformation $\vec{x} \mapsto A\vec{x}$ is not one-to-one. What else can you say about this transofrmaiton? Justify your answer.
    \question If $A$ is an $n \times n$ matrix and the transformation $x\mapsto A\vec{x}$ is one-to-one, what else can you say about this transformation? Justify your answer.
    \question Suppose $A$ is an $n \times n$ matrix with the property that the equation $A\vec{x}=\vec{b}$ has at least one solution for each $\vec{b}$ in $\mathbb{R}^n$. Without using Theorems 5 or 8, explain why each equation $A\vec{x} = \vec{b}$ has in fact one solution.
    \question Suppose $A$ is an $n \times n$ matrix with the property that the equation $A\vec{x}=\vec{0}$ has only the trivial solution. Without using the Invertible Matrix Theorem, explain directly why the equation $Aa\vec{x}=\vec{b}$ must have a solution for each $\vec{b}$ in $\mathbb{R}^n$.
   %=============
    \begin{center}\rule{4in}{0.1mm}\end{center}

    \uplevel{In Exercises 41 and 42, $T$ is a linear transformation from $\R{2}$} to $\R{2}$. Show that $T$ is invertible and find a formula for $T^{-1}$.
    \question $T(x_1,x_2) = (-9x_1 + 7x_2, 4x_1-3x_2)$
    \question $T(x_1,x_2) = (6x_1-8x_2,-5x_1+7x_2)$
    \question Let $T:\R{n}\rightarrow\R{n}$ be an invertible linear transformation. Explain why $T$ is both one-to-one and onto $\R{n}$. Use equations (1) and (2). Then give a second explanation using one or more theorems.
    \question Let $T$ be a linear transformation that msps $R{n}$ onto $\R{n}$. Show that $T^{-1}$ exists and maps $\R{n}$ onto $\R{n}$. Is $T^{-1}$ also one-to-one?
    \question Suppose $T$ and $U$ are linear transformations from $\R{n}$ to $\R{n}$ such that $T(U\vec{x}) = \vec{x}$ for all $\vec{x}$ in $\R{n}$. Is it true that $U(T\vec{x}) = \vec{x}$ for all $\vec{x}$ in $\R{n}$? Why or why not?
    \question Suppose a linear transofrmation $T:\R{n}\rightarrow\R{n}$ has the property that $T(\vec{u}) = T(\vec{varepsilon})$ for some pair of distinct vectors $\vec{u}$ and $\vec{v}$ in $\R{n}$. Can $T$ map $\R{n}$ onto $\R{n}$? Why or why not?
    \question Let $T:\R{n}\rightarrow\R{n}$ be an invertible linear transformation, and let $S$ and $U$ be functions from $\R{n}$ into $\R{n}$ such that $S(T(\vec{x})) = \vec{x}$ and $U(T(\vec{x})) = \vec{x}$ for all $\vec{x}$ in $\R{n}$. Show that $U(\vec{v}) = S(\vec{v})$ for all $\vec{v}$ in $\R{n}$. This will show that $T$ has a unique inverse, as asserted in Theorem 9. Hint: Given any $\vec{v}$ in $\R{n}$, we can write $\vec{v} = T(\vec{x})$ for some $\vec{x}$. Why? Compute $(\vec{v})$ and $U(\vec{v})$.
    \question Suppose $T$ and $S$ satisfy the invertibility equations (1) and (2), where $T$ is a linear transformation. Show directly that $S$ is a linear transformation. Hint: Given $\vec{u}$, $\vec{v}$ in $\R{n}$, let $\vec{x}=S(\vec{u})$, $\vec{y} = S(\vec{v})$. Then $T(\vec{x})=\vec{u}$, $T(\vec{y}) = \vec{v}$. Why? Apply $S$ to both sides of the equation $T(\vec{x}) + T(\vec{y}) = T(\vec{x}+\vec{y})$. Also, consider $T(c\vec{x}) = cT(\vec{x})$.

\end{questions} 
\end{document}